\documentclass[10pt,a4paper]{amsart}
\usepackage{amsmath,amssymb,amsthm,graphicx}
\usepackage{amsmath}
\usepackage{amssymb}
\usepackage{hyperref}
\newtheorem{Theorem}{Theorem}[section]
\newtheorem{Corollary}[Theorem]{Corollary}
\newtheorem{Lemma}[Theorem]{Lemma}
\newtheorem{Proposition}[Theorem]{Proposition}
{ \theoremstyle{definition}
	\newtheorem{Definition}[Theorem]{Definition}
	\newtheorem{Note}[Theorem]{Note}
	\newtheorem{Example}[Theorem]{Example}
	\newtheorem{Remark}[Theorem]{Remark} }

\usepackage[all]{xy}


\def\KK{\ensuremath\mathbb{K}}
\def\MM{\ensuremath\mathbb{M}}
\def\BB{\ensuremath\mathbb{B}}
\def\GL{\ensuremath\mathbf{GL}}
\def\UU{\ensuremath\mathbf{U}}
\def\PP{\ensuremath\mathbb{P}}
\def\CC{\ensuremath{\mathbb{C}}}
\def\NN{\ensuremath{\mathbb{N}}}
\def\QQ{\ensuremath{\mathbb{Q}}}
\def\RR{\ensuremath{\mathbb{R}}}
\def\ZZ{\ensuremath{\mathbb{Z}}}

\newcommand{\profinQ}{\mathfrak{Q}}
\newcommand{\reals}{\mathbb{R}}
\newcommand{\complexs}{\mathbb{C}}
\newcommand{\naturals}{\mathbb{N}}
\newcommand{\integers}{\mathbb{Z}}
\newcommand{\rationals}{\mathbb{Q}}
\newcommand{\quaternions}{\mathbb{H}}
\newcommand{\hyperbolic}{\mathbb{H}}
\newcommand{\K}{\mathbb{K}}
\newcommand{\realsorcomplexs}{\mathbb{K}}
\newcommand{\finitefield}{\mathbb{F}}
\newcommand{\Proj}{\mathbb{P}}
\DeclareMathOperator{\id}{id}
\DeclareMathOperator{\Ind}{Ind}
\newcommand{\boundary}[1]{\partial#1}
\newcommand{\boundedops}{\mathcal{B}}
\newcommand{\bundleover}{\!\!\downarrow\!\!}
\newcommand{\abs}[1]{\left\lvert#1\right\rvert} %absolute value
%\newcommand{\bigabs}[1]{\left\lvert#1\right\rvert} %absolute value
\newcommand{\norm}[1]{\left\lVert#1\right\rVert}
\newcommand{\generate}[1]{\langle#1\rangle}
\newcommand{\tensor}{\otimes}
\newcommand{\into}{\hookrightarrow}
\newcommand{\onto}{\twoheadrightarrow}
\newcommand{\iso}{\cong}
\newcommand{\disjointunion}{\amalg}
\newcommand{\subgroup}{\leq}
\newcommand{\supergroup}{\geq}
\newcommand{\superset}{\supset}
\newcommand{\normalsubgroup}{\lhd}
\newcommand{\semiProd}{\rtimes}
\newcommand{\semiprod}{\semiProd}
\newcommand{\dividesnot}{\nmid}
\newcommand{\divides}{\mid}
\DeclareMathOperator{\Pos}{Pos}
\DeclareMathOperator{\Aut}{Aut}
\DeclareMathOperator{\supp}{supp} %support
\DeclareMathOperator{\clos}{clos} %closure
\DeclareMathOperator{\im}{im} %image
\DeclareMathOperator{\vol}{vol} %Volume
\DeclareMathOperator{\diam}{diam} %Diameter
\DeclareMathOperator{\dist}{dist} %Distance
\DeclareMathOperator{\ord}{ord} %order
\DeclareMathOperator{\End}{End} %Endomorphisms
\DeclareMathOperator{\Hom}{Hom} %Homomorphisms
\DeclareMathOperator{\Span}{span}
\DeclareMathOperator{\spin}{spin}
\DeclareMathOperator{\diag}{diag}
\DeclareMathOperator{\spec}{Spec}
\DeclareMathOperator{\trace}{trace}
\DeclareMathOperator{\tr}{tr}
\DeclareMathOperator{\pr}{pr}
%\DeclareMathOperator{\res}{res} %Restriction
\DeclareMathOperator{\grad}{grad} %Restriction
% \DeclareMathOperator{\Sp}{Sp} %Spur % schon definiert???
\DeclareMathOperator{\fl}{f\/l} %flip
\DeclareMathOperator{\coker}{coker}
\DeclareMathOperator{\Char}{Char}
\DeclareMathOperator{\ind}{ind}
\DeclareMathOperator{\sgn}{sgn}
\DeclareMathOperator{\Ric}{Ric} % Ricci curvature
\DeclareMathOperator{\scal}{scal} % Scalar curvature
\DeclareMathOperator*{\dirlim}{dirlim}
\DeclareMathOperator*{\invlim}{invlim}
\DeclareMathOperator{\rank}{rank}
\DeclareMathOperator{\coind}{coind}
\DeclareMathOperator{\res}{res}
\newcommand{\bigdisjointunion}{\bigsqcup}
\DeclareMathOperator{\Orb}{Orb}
\DeclareMathOperator{\GTOP}{G-TOP}
\newcommand{\forget}[1]{}
\newenvironment{Forget}{}{}
\newcommand{\innerprod}[1]{\langle #1 \rangle}
\usepackage{cleveref}
\crefname{originaldefinition}{original definition}{original definitions}
\Crefname{originaldefinition}{Original definition}{Original definitions}
\makeatletter
\newlabel{def:Stolz_groups}{{3.1}{source}{}{sourcepointer-1}{}}
\newlabel{def:Stolz_groups@cref}{{[originaldefinition][0][]3.1}{[1][0][]source}{}{}{}}
\makeatother
\title{The equivariant Stolz positive scalar curvature exact sequence is exact}
\author{Massimiliano Puglisi, Thomas Schick, Vito Felice Zenobi}
\date{doi:10.3842/SIGMA.2025.093; CC BY 4.0}
\begin{document}
\maketitle
\noindent\textit{Editorial note.} All delivered mathematical text is the author's own, copied from the pinned publisher source sources/geometry-fl04-stolz-positive-scalar-curvature-st/pinned-original-source.tex; no statement or proof step was rewritten, reordered or removed.\par\medskip
\noindent\textit{Editorial note.} Preamble provenance: the retained preamble is the source's own layout adapter, if any, plus the author's macro block at lines 11-108 of sources/geometry-fl04-stolz-positive-scalar-curvature-st/pinned-original-source.tex (everything between the publisher class metadata and \textbackslash{}begin\{document\}, comments excluded). Nothing was added, no macro was invented, and the publisher document class is deliberately not reproduced; the excerpt serves only as preamble, never as a span coordinate.\par\medskip
\noindent\textit{Editorial note.} Metadata provenance: source\_id and source\_version were taken from sources/geometry-fl04-stolz-positive-scalar-curvature-st/license-observation.json; the title from SOURCE.json; the author list from the source metadata record; the source URL from the first non-empty of source\_url, source\_urls[0] and the evidence's fixed\_version\_url/license\_url. The licence evidence file used by this card is sources/geometry-fl04-stolz-positive-scalar-curvature-st/license-observation.json, recorded in license.evidence\_path so that the binding target is explicit.\par\medskip
\noindent\textit{Editorial note.} External references are navigation pointers only; every deliberately kept step of the author's argument stays inside the counted spans. The displayed numbers of those pointers are the author's own numbering under the author's \textbackslash{}newtheorem declarations (section.index as declared), computed mechanically by the lane's scanner.\par\medskip
\noindent\textit{Editorial note.} Bibliography provenance: the source stages its bibliography inside the pinned original (its bibliography block starts at source line 1361); the single entry cited inside the counted proof is reproduced verbatim in this task's own bibliography.tex, extracted from line 1417 of the pinned file and not edited in any way.\par\medskip
\noindent\textit{Editorial note.} The counted proof has one reference to another labelled item, the definition labelled def:Stolz\_groups, recorded as a navigation stub of kind definition.\par\medskip
\noindent\textit{Editorial note.} Assembly history: r001 failed with 'Text line contains an invalid character' because the per-card bibliography file had been written through a shell heredoc that turned the leading backslash of \textbackslash{}begin\{thebibliography\} into a backspace character, so the entry never registered and the citation stayed undefined; r002 rewrites that file with the file-writing tool instead of a shell script. The entry text itself was already correct and remains verbatim from line 1417 of the pinned source.\par\medskip
\noindent\textit{Allowed standard prerequisites (editorial).}\begin{itemize}
\item The three equivariant bordism groups appearing in the sequence, with their boundary, forgetful and closed-manifold maps. Same source, Definition labelled def:Stolz\_groups; invoked in the counted statement and used throughout the counted proof. Fixes the objects and the three maps whose exactness is asserted.
\item The classical non-equivariant Stolz exact sequence. Cited inside the counted proof (the long exact sequence (4.4) of Stolz); the author gives the equivariant argument in full. Locates the statement in the literature while the proof supplied is self-contained.
\end{itemize}\medskip
\ifdefined\hypertarget\hypertarget{sourcepointer-1}{}\fi
\noindent\textit{Source reference.} Original definition 3.1: sources/geometry-fl04-stolz-positive-scalar-curvature-st/pinned-original-source.tex, source lines 368--368. Navigation pointer for \textbackslash{}ref\{def:Stolz\_groups\} used inside the counted span (line 368); the argument that consumes it is retained in the span.\par\medskip
\begin{Proposition}\label{prop:Stolz_seq}
		The abelian groups defined in Definition~{\rm\ref{def:Stolz_groups}} fit into the
		following $G$-equivari\-ant version of the Stolz positive scalar curvature exact sequence:
		\[
			\cdots \to R^{\rm spin}_{n+1}(X)^{G} \to \mathrm{Pos}^{\rm spin}_{n}(X)^{G} \to \Omega^{\rm spin}_{n}(X)^{G} \to R^{\rm spin}_{n}(X)^{G} \to \cdots,
		\]
		where the first map sends a manifold to its boundary, the second one
		is the forgetful map $($i.e., it forgets the metric of positive scalar curvature$)$
		and the last one considers a closed manifold as a~manifold with empty boundary.
	\end{Proposition}

	\begin{proof}
		This is a rather direct consequence of the definitions and well known to the
		experts, at least non-equivariantly, stated, e.g., in \cite[long exact sequence~(4.4)]{Stolz} (but
		without proof). The
		argument for the
		$G$-equivariant case is exactly the same as the classical non-equivariant
		situation. As we are not aware of a treatment available in the journal
		literature, we follow the referee's suggestion and give a complete account
		of the argument here.
		
		First, that the composition of two consecutive maps is zero indeed is a
		direct consequence of the definitions:
		\begin{itemize}\itemsep=0pt
			\item Given $[W,f,g]\in R^{\rm spin}_{n+1}(X)^G$, its image in $\Omega_n(X)^G$ is
			represented by $(\boundary W,f|_{\boundary W})$ which indeed is a boundary, hence represents
			$0$.
			\item Given $[M,f,g]\in \mathrm{Pos}_n^{\rm spin}(X)^G$, its image in $R^{\rm spin}_n(X)^G$ is
			the manifold $M$ with empty boundary. It is bordant in $R^{\rm spin}_n(X)^G$
			to $\varnothing$ and hence represents $0$ via the
			bordism $M\times [0,1]$, where $(M\times \{1\},g)$ is a positive scalar curvature bordism of
			$\varnothing =\boundary M$ to $\varnothing=\boundary\varnothing$.
			\item Given $[M,f]\in \Omega^{\rm spin}_n(X)^G$, its image in
			$\mathrm{Pos}_{n-1}^{\rm spin}(X)^G$ is represented by $\boundary M=\varnothing$, hence
			represents $0$.
		\end{itemize}
		
		For the opposite inclusions of the kernels in the image, the constructions
		are almost as straight-forward:
		\begin{itemize}\itemsep=0pt
			\item Assume that $[M,f]\in \Omega^{\rm spin}_n(X)^G$ is mapped to $0$ in
			$R^{\rm spin}_n(X)^G$. This means that there is a null-bordism $(W,F)$. In the
			case at hand, $\boundary W$ has two disjoint parts: on the one hand~$(M,f)$
			and on the other hand $\bigl(M',f',g'\bigr)$ which is a bordism from $\boundary
			M=\varnothing$ to $\boundary\varnothing=\varnothing$ and which is equipped with
			a metric $g$ of positive scalar curvature. But this means, by definition,
			that $[M,f]=\big[M',f'\big]\in \Omega^{\rm spin}_n(X)^G$ where $M'$ is equipped with a
			metric $g$ of positive scalar curvature. Hence we have the preimage
			\smash{$\big[M',f',g\big]\in \mathrm{Pos}^{\rm spin}_n(X)^G$} of the initial equivariant bordism class
			$(M,f)$.
			\item Assume that $[M,f,g]\in \mathrm{Pos}^{\rm spin}_n(X)^G$ is mapped to $0$ in
			$\Omega^{\rm spin}_n(X)^G$. This means that there is a bordism $(W,F)$ with
			$\boundary(W,F)=(M,f)$. But then $[W,F,g]$ represents a class in~$R^{\rm spin}_{n+1}(X)^G$ which is mapped to $[M,f,g]\in \mathrm{Pos}^{\rm spin}_n(X)^G$.
			\item Finally, assume that \smash{$[W,f,g]\in R^{\rm spin}_{n+1}(X)^G$} is mapped to $0$
			in $\mathrm{Pos}_n^{\rm spin}(X)^G$. This means that there is a second bordism
			$\bigl(W',f',g'\bigr)$ whose boundary is $(\boundary W,f|_{\boundary W},g)$.
			We now consider the closed $G$-manifold $Z:= W\cup_{\boundary W} W'$ with
			map $F:= f\cup_{\boundary W} f'\colon Z\to X$ and the bordism $B:=Z\times
			[0,1]$ with map $F\circ \operatorname{pr}_Z\colon B\to X$. The boundary of $B$ consists
			of three parts: the internal boundary $W'$ between $\boundary W$ and
			$\varnothing =\boundary Z$, equipped with the metric~$g'$ of positive scalar
			curvature extending $g$. The other two boundary parts are $W$ and $Z$, and
			by definition $(B,F\circ \operatorname{pr}_Z, g')$ is a bordism in $R^{\rm spin}_n(X)^G$
			between $(W,f,g)$ and $(Z,f,\varnothing)$ and~${[Z,f]\in
			\Omega^{\rm spin}_n(X)^G}$ represents a preimage of $[W,f,g]\in R^{\rm spin}_n(X)^G$.\hfill $\qed$
		\end{itemize} \renewcommand{\qed}{}
	\end{proof}

% Bibliography excerpt for this card, copied verbatim from the pinned publisher source
% sources/geometry-fl04-stolz-positive-scalar-curvature-st/pinned-original-source.tex, whose own
% bibliography block starts at line 1361 and whose entry for the key cited in this card's proof is at
% line 1417. Only that entry is reproduced, character for character, with no reformatting or addition.
\begin{thebibliography}{99}

\bibitem{Stolz}
Stolz S., Concordance classes of positive scalar curvature metrics, {P}reprint, University of Notre Dame, 1999, \url{https://web.archive.org/web/20240430033820/https://www3.nd.edu/~stolz/concordance.ps}.

\end{thebibliography}
\end{document}
